\section{Introduction}
Transportation networks are a critical component of modern societies. They support economic activity, territorial cohesion, access to employment and services, and the movement of passengers and goods at multiple spatial and temporal scales. However, the goals of network design depend strongly on the application. In road networks, the emphasis is typically placed on reducing travel times and mitigating congestion, and the underlying optimization is usually framed as a network design problem on top of a user-equilibrium traffic-assignment model \cite{Sheffi1985,YangBell1998,MagnantiWong1984}. In rapid transit and public transport systems, planners are often more concerned with coverage, accessibility, and social cohesion, and the relevant models combine line planning, frequency setting and infrastructure deployment under public-service criteria \cite{GuihaireHao2008,LaporteMesaOrtegaPerea2011,RealRojas2026TRB}. In contrast, privately operated transportation systems are mainly designed from a profitability viewpoint, where the network configuration must balance service attractiveness against operating and infrastructure costs, and where market capture against incumbent competitors plays a central role \cite{AdlerSmilowitz2007,BelobabaOdoniBarnhart2015}. From an optimization perspective, these settings all lead to network design problems, but with markedly different objective functions and modeling choices.

Within this broad class of problems, hub-and-spoke networks play a particularly important role in air transportation \cite{OKelly1987,Campbell1994Integer,AdlerSmilowitz2007,BelobabaOdoniBarnhart2015}. Structurally, they can be interpreted as the union of several star-shaped subnetworks interconnected through a reduced set of hub airports. Spoke airports mainly feed or receive local traffic, whereas hubs concentrate transfer activity and allow the operator to connect many origin-destination markets with a comparatively limited number of direct links. For airlines, this architecture is attractive because it consolidates traffic, improves the use of aircraft and airport resources, and allows broad market coverage without operating every origin-destination pair as a direct service. Consequently, hub-and-spoke network design for a private airline is fundamentally a profit-driven problem in which the key decisions concern hub activation and the assignment of capacities across airports and routes. Passenger flow variables are incorporated not as direct operational decisions, but as a modeling device to represent how users may react to the offered network and, therefore, how much revenue each market can generate for the operator.

Classical formulations address these decisions through mixed-integer programming (MIP) models \cite{OKelly1987,Campbell1994Integer,AlumurKara2008,CampbellOKelly2012,FarahaniHekmatfarBolooriNikbakhsh2013}, since fixed-cost activation decisions are naturally represented with binary variables. This is the standard modeling approach for hub location and network design, but the resulting formulations become combinatorially challenging as the number of airports, markets, and operating restrictions grows, with complexity increasing very rapidly once multicommodity interactions and fixed costs are included \cite{AlumurKara2008,CampbellOKelly2012}. Modeling passenger behavior creates an additional challenge. Simplified all-or-nothing assignment rules are linear and easy to embed in optimization models \cite{Sheffi1985,MagnantiWong1984}, but they provide a poor representation of how users choose among competing itineraries. Logit-based demand models are considerably more realistic because they allow the captured share of each market to vary smoothly with the perceived utility of the available alternatives \cite{BenAkivaLerman1985,Train2009}. Yet, when introduced into a network design problem, they typically give rise to nonlinear and nonconvex constraints. In practice these constraints are handled either through piecewise-linear mixed-integer approximations of the logit relation \cite{Vielma2015}, which substantially inflate the number of binary variables and rapidly degrade the scalability of state-of-the-art branch-and-bound solvers, or through exponential-cone reformulations \cite{LubinYamangilBentVielma2018}, which preserve convexity but remain numerically demanding because they couple every market with cone constraints whose feasible set is hard to explore at scale. As a result, both approaches make medium- and large-size instances of the airline hub-and-spoke design problem computationally prohibitive in industrial settings.

This master's thesis builds on a recently proposed continuous bilevel optimization framework for transportation network design \cite{RealRojas2026TRB,CandesWakinBoyd2008}, in which reweighted $\ell_1$ sparsity is used to approximate the fixed costs associated with airport, hub, and route activation, while entropy-based regularization recovers logit market shares without explicitly imposing nonconvex share equations \cite{RealRojas2026TRB,Oppenheim1993}. On top of that continuous lower-level model we develop the main contribution of this work: we cast the airline design problem as a \emph{bilevel} optimization model. The upper level represents the airline as a profit-maximizing decision-maker that selects market-specific coefficients, whereas the lower level returns the continuous network design induced by those coefficients. In this way, the calibration of the model is performed jointly with the network design, instead of being delegated to an exogenous heuristic. To solve the resulting bilevel problem we adopt a penalty-based reformulation in the spirit of recent first-order bilevel algorithms \cite{AiyoshiShimizu1984,LuMei2024,Jiang2024Primal}, combined with an outer iteration that updates the reweighted $\ell_1$ terms, thereby improving the approximation of real fixed costs across successive solves.

The remainder of this master's thesis is organized as follows. We first review the relevant literature on hub-and-spoke network design, continuous approximations of fixed-cost models, demand modeling, and bilevel optimization. We then describe the airline problem setting and present the mixed-integer baseline together with the continuous lower-level formulation. Next, we introduce the bilevel reformulation and the associated penalty-based solution strategy. Finally, we report computational experience on synthetic Spanish networks of increasing size and on a realistic Madrid-centred network with Iberia-style international connections, and close with the main conclusions.
